\documentclass[10pt]{article}
\usepackage[a4paper, left=2.54cm, right=2.54cm, top=3cm, bottom=3cm]{geometry}
\usepackage[spanish]{babel} % Para fechas y formato en español
\usepackage{amsmath, amssymb} % Para notación matemática avanzada
\usepackage{multicol}
\usepackage{enumitem}
\usepackage{fancyhdr}
\usepackage{titlesec}

% Configuración de página
\pagestyle{fancy}
\fancyhf{}
\fancyhead[L]{Academia Prepolicial - Conjuntos}
\fancyhead[R]{\today}
\fancyfoot[C]{\thepage}
\renewcommand{\headrulewidth}{0.8pt}

% Configuración de títulos y espaciado
\titleformat{\section}[block]{\large\bfseries\centering}{\thesection}{1em}{}
\setlength{\columnsep}{1cm}
\setlength{\parskip}{0.8em}
\setlength{\parindent}{0em}

\begin{document}

	\begin{center}
		\Large\bfseries Conjuntos
	\end{center}

	\section*{Ejercicios}
	\begin{multicols}{2}
		\begin{enumerate}[label=\textbf{\arabic*.}, itemsep=0.4cm]
			\item Calcula la suma de los elementos del conjunto $A$ si:
			\[
			A = \{ (2y - 3) \in \mathbb{Z}, \, 2 \leq \sqrt{3y - 2} \leq 5 \}.
			\]

			\item Dado el conjunto:
			\[
			A = \{a; (b;c); d\},
			\]
			¿cuántos de los siguientes enunciados son incorrectos?
			\begin{enumerate}[label=\Roman*.]
				\item $\{b, c\} \subseteq A$
				\item $\{b, c\} \in A$
				\item $\{b, c\} \not\subseteq A$
				\item $c \in A$
				\item $\{a\} \subseteq A$
				\item $\{a\} \in A$
			\end{enumerate}

			\item Dados los conjuntos:
			\[
			A = \{1, 2, 3, 4, 5, 6\}, \quad B = \{2, 3, 4, 8, 3\},
			\]
			calcula el número de subconjuntos de $A$ más el número de subconjuntos de $B$.

			\item Si el subconjunto $A$ es unitario y es igual al conjunto $B$, calcula $a \times b$, donde:
			\[
			A = \{\sqrt{a+5};4\}, \quad B = \{\sqrt[3]{b + 7^2}\} .
			\]

			\item Si los conjuntos $A$ y $B$ son unitarios, calcula $x^y$, donde:
			\[
			A = \{\sqrt[3]{x+3};2\}, B = \{y^x;32\}.
			\]

			\item Calcula el cardinal de:
			\[
			C = \{ (2x+3) \in \mathbb{N} \mid 3\leq x <  5 \}.
			\]

			\item Calcula el cardinal del conjunto $B$ si:
			\[
			B = \{ x^2 + 2 \mid x \in \mathbb{Z} \mid -2 \leq x < 3 \}.
			\]


			\item Calcula $n(A) + n(B)$ si:
			\[
			A = \{ \frac{x+3}{2} \in \mathbb{N} \mid 2 \leq x \leq 5 \}, \quad B = \left\{ \frac{y+4}{3} \mid y = A.
			\]

			\item Dados los conjuntos:
			\[
			A = \{ x \in \mathbb{N} \mid 25x^2 + 5x + 10 + 1 = 0 \}, \quad
			\]
			\[
			C = \{ \frac{1}{x} \in \mathbb{R} \mid 4x^2 - 4x + 1<=  0 \}
			\]
			calcula $n(A) + n(C)$.

			\item El número de subconjuntos de un conjunto de
			n + 2 elementos excede al doble del número de
			subconjuntos de un conjunto de n – 2 elementos
			en 224. Calcula el valor de “n”.
		\end{enumerate}
	\end{multicols}

\end{document}